% Sections 9.2.5 and 9.2.6, short version. Labels: fig:qual-cs-1..3, tab:paths-cs-1..3 (fig_examples_cs.tex),
% fig:graph-cs-2 (fig_graph_cs_creativity.tex), fig:qual-example-ot, tab:path-interpretations
% (fig_example_biology_ot.tex). sec:ablation, sec:faith, \cite{nbfnet}, \cite{gfmrag} are placeholders.

\subsection{Qualitative Analysis}
\label{sec:qualitative}

A cross-domain inspiration shares a mechanism with the query, not a vocabulary, so dense retrieval
misses it. We trace what each component of SciGraphIR adds on such queries by reporting the rank of
the gold at every stage of the cumulative ablation: Qwen3 cosine, multi-view scorer, OpenIE entity
graph, SciAffordGraph without CCMP, and the full model. Examples were selected by the ladder itself:
of the 926 cross-field golds in SIR-4 CS, we kept those whose rank improves at every stage.

\paragraph{Computer science to psychology (Figure~\ref{fig:qual-cs-2}).}
The query asks why language models generate less diverse ideas than people. The gold is a 1993
cognitive-psychology experiment: exposure to examples constrains the ideas people generate. The
query never mentions psychology; cosine ranks the paper 169th and no dense baseline puts it in the
top 25. The scorer's hypothetical answers bring it to 49, the entity graph cannot reach it (48), the
frame graph reaches it at 25, CCMP at 18; the graph channel alone ranks it 2nd.
Figure~\ref{fig:graph-cs-2} shows the route: the query's \textsc{function} frame on incubation leads
to a psychology paper on fixation, its \textsc{domain} node leads to the gold, and CCMP amplifies
that crossing (gate 1.15). This is the argument a researcher would make: idea convergence is a
fixation effect, and fixation is studied in cognitive psychology.

\paragraph{Biology to computer vision (Figure~\ref{fig:qual-example-ot}).}
A protein--RNA binding-site query is inspired by a vision paper on class-imbalanced recognition.
The hypothetical answer names the function the query needs, handling class imbalance; the frame graph
links the query's \textsc{limitation} frame in two hops to the optimal-transport method the vision
paper contributes. Cosine 113, scorer 13, graph 3.

\paragraph{Which graph.}
The OpenIE entity graph, built from the same corpus, reaches papers only through co-mention. It has
no route to the psychology paper and ranks the vision paper 3{,}329th against the frame graph's 3
(Table~\ref{tab:path-interpretations}). Entity graphs connect papers that share words; the frame
graph connects a need in the query to the method, finding, or limitation that answers it.

\paragraph{What CCMP does.}
With the gate switched off on the same weights, the routes are unchanged (14 of 15 top routes
identical, grey weights in Tables~\ref{tab:paths-cs-1} to \ref{tab:paths-cs-3}). CCMP changes their
weights: it raises specific method routes over generic task routes (Table~\ref{tab:paths-cs-1}) and
moves the psychology gold from 4 to 2 in the graph channel. Its effect is one to two ranks per query.

\subsection{Path Interpretations}
\label{sec:paths}

Following NBFNet~\cite{nbfnet} and GFM-RAG~\cite{gfmrag}, the importance of a path $P$ to the
graph-channel score $g_q(d)$ of document $d$ is the sum of the partial derivatives of the score with
respect to the weights of its edges at the layers they were used,
\begin{equation}
  P_1, \dots, P_k \;=\; \operatorname*{arg\,top\text{-}k}_{P}\;
  \sum_{(h, r, t)\in P} \frac{\partial\, g_q(d)}{\partial\, w^{(l)}_{(h,r,t)}},
  \label{eq:path-interp}
\end{equation}
with the top-$k$ paths found by beam search from the query's seed frames to $d$. We additionally
record the CCMP gate applied to each hop's sender, normalised by the frontier mean, so a value above
1 marks an amplified hop. Tables~\ref{tab:paths-cs-1} to \ref{tab:paths-cs-3} and
\ref{tab:path-interpretations} list the top routes for the examples above.

The routes are short and typed: one to four hops from a \textsc{function}, \textsc{limitation}, or
\textsc{method} frame of the query to the frame the gold contributes, achieves, or overcomes.
Cross-domain routes cross at a bridge paper or a domain node, not at shared words. And the routes
are the same with and without CCMP; the gate sets their weight. Section~\ref{sec:faith} tests these
three properties over all cross-field queries.
